\documentclass[11pt]{article}
\usepackage[margin=1in]{geometry}
\usepackage[T1]{fontenc}
\usepackage{lmodern}
\usepackage{amsmath}
\usepackage{amssymb}
\usepackage[ruled,vlined,linesnumbered]{algorithm2e}
\usepackage[colorlinks=true,linkcolor=blue,citecolor=blue,urlcolor=blue]{hyperref}
\usepackage[nameinlink,noabbrev]{cleveref}

\SetKwInput{KwData}{Data}
\SetKwInput{KwResult}{Result}
\SetKwComment{Comment}{$\triangleright$\ }{}
\crefname{algocf}{algorithm}{algorithms}
\Crefname{algocf}{Algorithm}{Algorithms}

\title{Algorithm Float Conversion Stress Paper}
\author{Stage Two Corpus}
\date{}

\begin{document}
\maketitle

\section{Overview}

Algorithm environments often combine captions, line numbers, mathematical symbols, nested control flow, and cross-references. This fixture places pseudocode in a ruled float so that a converter must preserve both the program-like layout and the surrounding scholarly reference structure.

The screening procedure in \cref{alg:adaptive-screen} computes a risk score for each candidate, updates a calibration term, and records the first threshold crossing. The surrounding prose uses the algorithm as a numbered scholarly object rather than as plain code.

\begin{algorithm}[tbp]
\DontPrintSemicolon
\caption{Adaptive screen with calibrated stopping}
\label{alg:adaptive-screen}
\KwData{Candidate set $\mathcal{J}$, features $x_j$, threshold $\tau$, and maximum rounds $T$}
\KwResult{Accepted set $A$ and audit log $L$}
$A \leftarrow \varnothing$; $L \leftarrow [\,]$; $\lambda \leftarrow 0$\;
\For{$t \leftarrow 1$ \KwTo $T$}{
  $B_t \leftarrow \varnothing$\;
  \ForEach{$j \in \mathcal{J} \setminus A$}{
    $s_j \leftarrow w_t^{\top}x_j - \lambda$\;
    \uIf{$s_j \geq \tau$}{
      $A \leftarrow A \cup \{j\}$\;
      append $(t,j,s_j,\mathrm{accept})$ to $L$\;
    }
    \ElseIf{$\tau - s_j \leq 0.05$}{
      $B_t \leftarrow B_t \cup \{j\}$\;
      append $(t,j,s_j,\mathrm{borderline})$ to $L$\;
    }
  }
  $\lambda \leftarrow \lambda + |B_t|/(1+|\mathcal{J}|)$\Comment*[r]{calibrate future rounds}
  \If{$B_t=\varnothing$ \KwSty{and} $A \neq \varnothing$}{
    \KwRet{$(A,L)$}\;
  }
}
\KwRet{$(A,L)$}\;
\end{algorithm}

\section{Interpretation}

The most important structural feature of \cref{alg:adaptive-screen} is the coexistence of mathematical notation and algorithmic indentation. When the float is converted, the line-number gutter, comments, and return statements need stable alignment because readers use them to distinguish the scoring rule from the stopping condition.

\end{document}
